% !TEX root = ../main.tex
\begin{figure}[H]
  \centering
  \resizebox{\textwidth}{!}{%
    \begin{tikzpicture}[
      papel/.style={nodoDiag, minimum height=0.9cm, font=\small},
      papelFin/.style={papel, draw=diagFocoBorde, fill=diagFocoFondo, very thick},
      lineaVida/.style={densely dotted, thick, draw=diagNodoBorde},
      mensaje/.style={-{Latex[length=2mm]}, semithick},
      devuelve/.style={-{Latex[length=2mm]}, semithick, densely dashed},
      devuelveFin/.style={-{Latex[length=2mm]}, very thick, draw=diagFocoBorde},
      etiqueta/.style={font=\small, midway, above, inner sep=1.5pt, fill=white},
      propia/.style={font=\small, anchor=west, inner sep=1.5pt, fill=white},
      bloque/.style={draw=diagGrupoBorde, semithick, densely dashed,
        fill=diagGrupo, fill opacity=0.4, text opacity=1,
        rounded corners=2pt},
      rotuloBloque/.style={font=\small\bfseries, anchor=north west,
        inner sep=2.5pt, text=diagGrupoBorde!65!black, fill=white},
      apunte/.style={font=\small, align=center, draw=diagNodoBorde,
        fill=white, inner sep=3pt, rounded corners=2pt},
      ]

      \node[papel]    at (0,-1.3)    {El estudiante};
      \node[papel]    at (3.7,-1.3)  {La conversación};
      \node[papel]    at (7.4,-1.3)  {El recuperador};
      \node[papel]    at (11.1,-1.3) {El generador};
      \node[papelFin] at (14.8,-1.3) {El modelo local};
      \foreach \x in {0, 3.7, 7.4, 11.1, 14.8}{
        \draw[lineaVida] (\x,-1.75) -- (\x,-16.4);
      }

      \draw[mensaje] (0,-2.3) -- node[etiqueta]{pregunta} (3.7,-2.3);
      \draw[mensaje] (3.7,-2.9) -- ++(0.8,0) -- ++(0,-0.45) -- ++(-0.8,0);
      \node[propia] at (4.6,-3.15) {resolver el ámbito y la anáfora};

      \node[bloque, fit={(-1.4,-3.9) (16.2,-5.3)}] (b1) {};
      \node[rotuloBloque] at (b1.north west) {el mensaje no pide un dato};
      \draw[devuelve] (3.7,-4.9) -- node[etiqueta]{respuesta fija} (0,-4.9);
      \node[apunte] at (10.6,-4.8)
      {un saludo, una despedida o preguntar por otro centro};

      \draw[mensaje] (3.7,-6.0) -- node[etiqueta]{consulta} (7.4,-6.0);
      \draw[mensaje] (7.4,-6.6) -- ++(0.8,0) -- ++(0,-0.45) -- ++(-0.8,0);
      \node[propia] at (8.3,-6.85) {incrustar y buscar los veinte más próximos};
      \draw[mensaje] (7.4,-7.6) -- ++(0.8,0) -- ++(0,-0.45) -- ++(-0.8,0);
      \node[propia] at (8.3,-7.85) {acotar: banda, corte relativo y suelo};
      \draw[devuelve] (7.4,-8.6) -- node[etiqueta]{fragmentos} (3.7,-8.6);

      \node[bloque, fit={(-1.4,-9.2) (16.2,-10.6)}] (b2) {};
      \node[rotuloBloque] at (b2.north west) {no queda ningún fragmento};
      \draw[devuelve] (3.7,-10.2) -- node[etiqueta]{respuesta fija} (0,-10.2);
      \node[apunte] at (10.6,-10.1)
      {sin contexto no se llama al modelo:\\es el estado en el que inventa};

      \draw[mensaje] (3.7,-11.3)
      -- node[etiqueta]{pregunta, fragmentos y turnos} (11.1,-11.3);
      \draw[mensaje] (11.1,-11.9) -- ++(0.8,0) -- ++(0,-0.45) -- ++(-0.8,0);
      \node[propia] at (12.0,-12.15) {armar el \textit{prompt}};
      \draw[mensaje] (11.1,-12.9) -- node[etiqueta]{\textit{prompt}} (14.8,-12.9);
      \draw[devuelve] (14.8,-13.5) -- node[etiqueta]{respuesta} (11.1,-13.5);
      \draw[mensaje] (11.1,-14.1) -- ++(0.8,0) -- ++(0,-0.45) -- ++(-0.8,0);
      \node[propia] at (12.0,-14.35)
      {comprobar cada titulación contra el catálogo};

      \node[bloque, fit={(-1.4,-15.0) (16.2,-16.4)}] (b3) {};
      \node[rotuloBloque] at (b3.north west) {nombra alguna que no existe};
      \draw[devuelve] (11.1,-16.0)
      -- node[etiqueta]{respuesta retirada entera} (0,-16.0);

      \draw[lineaVida] (0,-16.4) -- (0,-17.5);
      \draw[lineaVida] (11.1,-16.4) -- (11.1,-17.5);
      \draw[devuelveFin] (11.1,-17.1) -- node[etiqueta]{respuesta} (0,-17.1);
    \end{tikzpicture}%
  }
  \caption[Recorrido completo de una consulta]{Recorrido
      completo de una consulta. Los tres bloques son las salidas anticipadas, y
      las dos primeras resuelven la respuesta sin invocar al modelo generativo.
      Fuente:
      Elaboración propia.}
  \label{fig:secuencia_rag}
\end{figure}
